\section{Z}

% zahlen (to pay)
\begin{verbentry}{
  style = {default},
  toc = {zahlen (to pay)},
  name = {zahlen},
  english = {to pay},
  type = {regular},
  auxiliary = {haben}
}


\vspace{5pt}
\hrule
\vspace{5pt}

\begin{tabular}{rl}
\multicolumn{2}{l}{\textbf{PRÄSENS}}\\
ich & \textbf{zahle}\\
du & \textbf{zahlst}\\
er/sie/es & \textbf{zahlt}\\
wir & \textbf{zahlen}\\
ihr & \textbf{zahlt}\\
sie/Sie & \textbf{zahlen}\\
\end{tabular}
\hfill
\begin{tabular}{rl}
\multicolumn{2}{l}{\textbf{PRÄTERITUM}}\\
ich & \textbf{zahlte}\\
du & \textbf{zahltest}\\
er/sie/es & \textbf{zahlte}\\
wir & \textbf{zahlten}\\
ihr & \textbf{zahltet}\\
sie/Sie & \textbf{zahlten}\\
\end{tabular}
\hfill
\begin{tabular}{rl}
\multicolumn{2}{l}{\textbf{IMPERATIV}}\\
& \textbf{zahle (du)!}\\
& \textbf{zahlen wir!}\\
& \textbf{zahlt (ihr)!}\\
& \textbf{zahlen Sie!}\\
\end{tabular}

\vspace{10pt}

\begin{tabular}{rl}
\multicolumn{2}{l}{\textbf{PERFEKT}}\\
& haben + \textbf{gezahlt}\\
\end{tabular}
\hfill
\begin{tabular}{rl}
\multicolumn{2}{l}{\textbf{FUTURE I}}\\
& werden + \textbf{zahlen}\\
\end{tabular}
\hfill
\begin{tabular}{rl}
\multicolumn{2}{l}{\textbf{PLUSQUAMPERFEKT}}\\
& hatten + \textbf{gezahlt}\\
\end{tabular}

\vspace{10pt}

\begin{tabular}{lll}
\multicolumn{3}{l}{\textbf{MIT PRÄPOSITIONEN}}\\
& \textbf{zahlen für} + Akk & (to pay for)\\
& \textbf{zahlen an} + Akk & (to pay to)\\
\end{tabular}
\hfill
\begin{tabular}{lll}
\multicolumn{3}{l}{\textbf{TRENNBARE VERBEN}}\\
\vspace{20pt}
\end{tabular}
\vspace{-5pt}
\end{verbentry}

% zeigen (to show / point out)
\begin{verbentry}{
  style = {acc},
  toc = {zeigen (to show / point out)},
  name = {zeigen},
  english = {to show / point out},
  type = {regular},
  auxiliary = {haben}
}
 
    
     \vspace{5pt}

    \hrule
    
    \vspace{5pt}

  \begin{tabular}{rl}
     \multicolumn{2}{l}{\textbf{PRÄSENS} }\\
    ich & \textbf{zeige}\\
    du & \textbf{zeigst}\\
    er/sie/es & \textbf{zeigt}\\
    wir & \textbf{zeigen}\\
    ihr & \textbf{zeigt}\\
    sie/Sie & \textbf{zeigen}\\
  \end{tabular}
  \hfill
  \begin{tabular}{rl}
    \multicolumn{2}{l}{\textbf{PRÄTERITUM} }\\
    ich & \textbf{zeigte}\\
    du & \textbf{zeigtest}\\
    er/sie/es & \textbf{zeigte}\\
    wir & \textbf{zeigten}\\
    ihr & \textbf{zeigtet}\\
    sie/Sie & \textbf{zeigten}\\
  \end{tabular}
  \hfill
  \begin{tabular}{rl}
     \multicolumn{2}{l}{\textbf{IMPERATIV} }\\
   & \textbf{zeige (du)!}\\
   & \textbf{zeigen wir!}\\
   & \textbf{zeigt (ihr)!}\\
   & \textbf{zeigen Sie!}\\
   & \vspace{10pt}
  \end{tabular}

 \vspace{10pt}

 \begin{tabular}{rl}
     \multicolumn{2}{l}{\textbf{PERFEKT}}\\
   & haben + \textbf{gezeigt}\\
  \end{tabular}
\hfill
\begin{tabular}{rl}
    \multicolumn{2}{l}{\textbf{FUTURE I}}\\
   & werden + \textbf{zeigen}\\
  \end{tabular}
\hfill
\begin{tabular}{rl}
    \multicolumn{2}{l}{\textbf{PLUSQUAMPERFEKT}}\\
   & hatten + \textbf{gezeigt}\\
  \end{tabular}

 \vspace{10pt}

\begin{tabular}{lll}
      \multicolumn{3}{l}{\textbf{MIT PRÄPOSITIONEN}}\\
& \textbf{zeigen auf} + Akk & (to point to / indicate)\\
& \textbf{zeigen für} + Akk & (to show for / demonstrate)\\
& \textbf{zeigen mit} + Dat & (to show with)\\
& \textbf{zeigen von} + Dat & (to show of / about)\\
\end{tabular}
  \hfill
\begin{tabular}{lll}
      \multicolumn{3}{l}{\textbf{TRENNBARE VERBEN}}\\
 & \textbf{vor$|$zeigen} & (to demonstrate / exhibit)\\
 & \textbf{auf$|$zeigen} & (to point out / indicate)\\
 & \textbf{nach$|$zeigen} & (to show afterward / demonstrate later)\\
\vspace{10pt}
  \end{tabular}
  
  \vspace{-5pt}
\end{verbentry}

% zerstören (to destroy)
\begin{verbentry}{
  style = {default},
  toc = {zerstören (to destroy)},
  name = {zerstören},
  english = {to destroy},
  type = {regular, inseparable},
  auxiliary = {haben}
}
 
    
     \vspace{5pt}

    \hrule
    
    \vspace{5pt}

  \begin{tabular}{rl}
     \multicolumn{2}{l}{\textbf{PRÄSENS} }\\
    ich & \textbf{zerstöre}\\
    du & \textbf{zerstörst}\\
    er/sie/es & \textbf{zerstört}\\
    wir & \textbf{zerstören}\\
    ihr & \textbf{zerstört}\\
    sie/Sie & \textbf{zerstören}\\
  \end{tabular}
  \hfill
     \begin{tabular}{rl}
    \multicolumn{2}{l}{\textbf{PRÄTERITUM} }\\
    ich & \textbf{zerstörte}\\
    du & \textbf{zerstörtest}\\
    er/sie/es & \textbf{zerstörte}\\
    wir & \textbf{zerstörten}\\
    ihr & \textbf{zerstörtet}\\
    sie/Sie & \textbf{zerstörten}\\
  \end{tabular}
  \hfill
  \begin{tabular}{rl}
     \multicolumn{2}{l}{\textbf{IMPERATIV} }\\
   & \textbf{zerstöre (du)!}\\
   & \textbf{zerstören wir!}\\
   & \textbf{zerstört (ihr)!}\\
   & \textbf{zerstören Sie!}\\
   & \vspace{10pt}
  \end{tabular}

    \vspace{10pt}

 \begin{tabular}{rl}
     \multicolumn{2}{l}{\textbf{PERFEKT}}\\
   & haben + \textbf{zerstört}\\
  \end{tabular}
\hfill
   \begin{tabular}{rl}
    \multicolumn{2}{l}{\textbf{FUTURE I} }\\
   & werden + \textbf{zerstören}\\
  \end{tabular}
\hfill
   \begin{tabular}{rl}
      \multicolumn{2}{l}{\textbf{PLUSQUAMPERFEKT}}\\
   & hatten + \textbf{zerstört}\\
  \end{tabular}

   \vspace{10pt}

   \begin{tabular}{lll}
      \multicolumn{3}{l}{\textbf{MIT PRÄPOSITIONEN}}\\
& \textbf{zerstören durch} + Akk & (to destroy through / by)\\
& \textbf{zerstören mit} + Dat & (to destroy with)\\
\end{tabular}
  \hfill
   \begin{tabular}{lll}
\multicolumn{3}{l}{\textbf{TRENNBARE VERBEN}}\\
& \textbf{---} & \\
\end{tabular}

  \vspace{-5pt}
\end{verbentry}

% ziehen (to pull)
\begin{verbentry}{
  style = {default},
  toc = {ziehen (to pull)},
  name = {ziehen},
  english = {to pull},
  type = {irregular},
  auxiliary = {haben}
}
 
    
     \vspace{5pt}

    \hrule
    
    \vspace{5pt}

  \begin{tabular}{rl}
     \multicolumn{2}{l}{\textbf{PRÄSENS} }\\
    ich & \textbf{ziehe}\\
    du & \textbf{ziehst}\\
    er/sie/es & \textbf{zieht}\\
    wir & \textbf{ziehen}\\
    ihr & \textbf{zieht}\\
    sie/Sie & \textbf{ziehen}\\
  \end{tabular}
  \hfill
     \begin{tabular}{rl}
    \multicolumn{2}{l}{\textbf{PRÄTERITUM} }\\
    ich & \textbf{zog}\\
    du & \textbf{zogst}\\
    er/sie/es & \textbf{zog}\\
    wir & \textbf{zogen}\\
    ihr & \textbf{zogt}\\
    sie/Sie & \textbf{zogen}\\
  \end{tabular}
  \hfill
  \begin{tabular}{rl}
     \multicolumn{2}{l}{\textbf{IMPERATIV} }\\
   & \textbf{zieh(e) (du)!}\\
   & \textbf{ziehen wir!}\\
   & \textbf{zieht (ihr)!}\\
   & \textbf{ziehen Sie!}\\
   & \vspace{10pt}
  \end{tabular}

    \vspace{10pt}

 \begin{tabular}{rl}
     \multicolumn{2}{l}{\textbf{PERFEKT}}\\
  &  haben + \textbf{gezogen}\\
  \end{tabular}
\hfill
   \begin{tabular}{rl}
   
    \multicolumn{2}{l}{\textbf{FUTURE I} }\\
  &  werden + \textbf{ziehen}\\
 
  \end{tabular}
\hfill
   \begin{tabular}{rl}
      \multicolumn{2}{l}{\textbf{PLUSQUAMPERFEKT}}\\
  &  hatten + \textbf{gezogen}
  \end{tabular}
  


  \vspace{10pt}

   \begin{tabular}{lll}
      \multicolumn{3}{l}{\textbf{MIT PRÄPOSITIONEN}}\\
& \textbf{ziehen an} + Dat & (to pull on / at)\\
& \textbf{ziehen nach} + Dat & (to move to / after)\\
& \textbf{ziehen in} + Akk & (to move into)\\
& \textbf{ziehen aus} + Dat & (to move out of)\\
& \textbf{ziehen für} + Akk & (to pull for / draw for)\\
\vspace{10pt}
  \end{tabular}
  \hfill
   \begin{tabular}{lll}
      \multicolumn{3}{l}{\textbf{TRENNBARE VERBEN}}\\
 & \textbf{ein$|$ziehen} & (to move in / enlist / draw in)\\
 & \textbf{aus$|$ziehen} & (to move out / undress / extract)\\
 & \textbf{um$|$ziehen} & (to move / relocate / change clothes)\\
 & \textbf{zurück$|$ziehen} & (to pull back / withdraw)\\
 & \textbf{ab$|$ziehen} & (to subtract / withdraw / remove)\\
 & \textbf{auf$|$ziehen} & (to wind up / draw up / open)\\
 & \textbf{durch$|$ziehen} & (to pull through / carry out / complete)
  \end{tabular}


  \vspace{-5pt}
\end{verbentry}
